% ══════════════════════════════════════════════════════════════════════════════
%  EXPERIMENT 2 — MEASURING MASS AND VOLUME                            (Lesson 3)
% ══════════════════════════════════════════════════════════════════════════════
\experimentdiv{Measuring Mass and Volume}{Using an electronic balance \textbullet{} Reading a measuring cylinder \textbullet{} Precision}{3}\label{sec:2}
\lessoncodes{VC2S8I02 \textbullet{} VC2S8I03 \textbullet{} VC2S8I04 \textbullet{} VC2S8I05 \textbullet{} VC2S8I06}

\vocabhead
\begin{ctab}
\begin{tabular}{|SW{4.2cm}|W{11.6cm}|}
  \hline\hd & \thead{Word} & \thead{Definition (copy from the slide)} \\ \hline
  & \scaf{\textbf{Mass}} & \scaf{The amount of matter in an object, measured in grams (g) or kilograms (kg).} \\ \hline
  & \textbf{Volume} & \\ \hline
  & \textbf{Meniscus} & \\ \hline
  & \textbf{Tare} & \\ \hline
  & \textbf{Precision} & \\ \hline
\end{tabular}
\end{ctab}

\begin{lipanel}
\inquiryli{Use equipment to collect precise measurements and record them in a table with units.}
\understandli{Explain the difference between mass and volume.}
\sctitle
\begin{itemize}[leftmargin=1.4em, itemsep=0pt, topsep=1pt]
  \item I can tare an electronic balance and record a mass with its unit.
  \item I can read a measuring cylinder at eye level, from the bottom of the meniscus.
  \item I can record measurements in a table with the units in the column headings.
  \item I can explain the difference between mass and volume.
\end{itemize}
\end{lipanel}

\begin{skillsfocus}
\focus{PC4}{Read each scale to its smallest division, and choose the cylinder with the finest scale.}
\focus{PC6}{Measure each object three times and record every reading with its unit.}
\focus{PC7}{Use an electronic balance and a measuring cylinder correctly.}
\focus{PA1}{Record your data in tables, with units in the headings.}
\focus{PA3}{Calculate an average and a difference from your data.}
\end{skillsfocus}

% ─── TECHNIQUE ────────────────────────────────────────────────────────────────
\exppart{Technique spotlight}{Blue}

\begin{technique}{Using an electronic balance}{PC7}
\begin{minipage}[c]{0.6\linewidth}
\begin{enumerate}[leftmargin=1.6em, itemsep=1pt, topsep=0pt]
  \item Put the balance on a flat, steady bench, away from water.
  \item Press \textbf{TARE} with nothing on the pan, so it reads 0.00\,g.
  \item Place the object gently in the middle of the pan.
  \item Wait until the number stops changing.
  \item Record \textbf{every} digit shown, and the unit.
\end{enumerate}
\end{minipage}\hfill
\begin{minipage}[c]{0.36\linewidth}\centering
\balancefig{24.68\,g}\\[2pt]
{\footnotesize Record this as \textbf{24.68\,g}, not 24.7\,g.}
\end{minipage}

\medskip
\textbf{Taring with a container.} To find the mass of a liquid, put the empty beaker on the pan and press \textbf{TARE}. The balance now reads 0.00\,g, so when you pour in the liquid it shows the mass of the liquid alone.
\end{technique}

\begin{workedex}{The mass of some water without taring}
A beaker has a mass of 52.30\,g. With water in it, the balance reads 87.65\,g. Find the mass of the water.

\smallskip
\textbf{Working:} mass of water $=$ (beaker $+$ water) $-$ beaker $= 87.65 - 52.30 = 35.35$\,g.
\end{workedex}

\begin{wedo}{The mass of some water without taring}
A beaker has a mass of 48.20\,g. With water in it, the balance reads 83.65\,g. Find the mass of the water.
\workbox{3cm}
\end{wedo}

\begin{technique}{Reading a measuring cylinder}{PC4}
\begin{minipage}[c]{0.64\linewidth}
\begin{enumerate}[leftmargin=1.6em, itemsep=1pt, topsep=0pt]
  \item Stand the cylinder on a flat bench. Do not hold it up.
  \item \textbf{Work out the scale.} Take two labelled marks, find the difference between them, and divide by the number of spaces between them. That is the value of one small division.
  \item Bend down so your eye is \textbf{level} with the top of the water. Looking from above or below gives the wrong reading (\emph{parallax error}).
  \item Water curves down in the middle. This curve is the \textbf{meniscus}. Read the scale at the \textbf{bottom} of the meniscus.
  \item Record the reading and the unit, mL.
\end{enumerate}
\end{minipage}\hfill
\begin{minipage}[c]{0.32\linewidth}\centering
\begin{tikzpicture}
  \fill[Sky!22] (0,0) -- (0,1.36) .. controls (0.6,1.0) and (1.6,1.0) .. (2.2,1.36) -- (2.2,0) -- cycle;
  \draw[Ink, line width=0.9pt] (0,1.36) .. controls (0.6,1.0) and (1.6,1.0) .. (2.2,1.36);
  \draw[eq] (0,0) -- (0,2.2) (2.2,0) -- (2.2,2.2);
  \draw[OrangeDark, line width=0.8pt, dashed] (-0.3,1.09) -- (2.5,1.09);
  \node[font=\scriptsize, align=center, anchor=north] at (1.1,-0.05) {read the scale\\at the bottom\\of the meniscus};
  \node[font=\scriptsize] at (3.05,1.09) {eye};
  \fill[Ink] (2.75,1.09) circle (1.2pt);
\end{tikzpicture}
\end{minipage}
\end{technique}

\begin{workedex}{Reading a measuring cylinder}
\begin{minipage}[c]{0.64\linewidth}
Read the volume of water in this measuring cylinder.

\smallskip
\textbf{Scale:} the labelled marks 30 and 40 are 10\,mL apart, with 10 spaces between them. One small division is $10 \div 10 = 1$\,mL.

\smallskip
\textbf{Reading:} the bottom of the meniscus is 4 small divisions above 30, so the volume is $30 + 4 \times 1 = $ \textbf{34\,mL}.
\end{minipage}\hfill
\begin{minipage}[c]{0.3\linewidth}\centering
\readcyl{20}{50}{10}{10}{34}
\end{minipage}
\end{workedex}

\begin{wedo}{Reading a measuring cylinder}
\begin{minipage}[c]{0.64\linewidth}
\pt{a}What is the value of one small division?
\workbox{2.2cm}
\pt{b}Read the volume of water.
\answerbox{1cm}
\end{minipage}\hfill
\begin{minipage}[c]{0.3\linewidth}\centering
\readcyl{0}{25}{5}{5}{17}
\end{minipage}
\end{wedo}

\begin{technique}{Choosing the most precise equipment}{PC4}
The smaller the value of one small division, the more \textbf{precise} the measurement. Choose the \textbf{smallest} measuring cylinder that holds the whole volume. A 10\,mL cylinder can be read to 0.2\,mL, but a 250\,mL cylinder only to 10\,mL.
\end{technique}

\begin{workedex}{Choosing a measuring cylinder}
You need 8\,mL of water. You have a 10\,mL, a 50\,mL and a 100\,mL measuring cylinder.

\smallskip
\textbf{Choice:} the \textbf{10\,mL} cylinder. It holds 8\,mL and has the smallest divisions (0.2\,mL), so the reading is the most precise. In the 100\,mL cylinder, 8\,mL barely covers the base.
\end{workedex}

\begin{wedo}{Choosing a measuring cylinder}
You need 45\,mL of water. You have a 10\,mL, a 50\,mL and a 250\,mL measuring cylinder. Choose one and explain your choice.
\answerbox{5cm}
\end{wedo}

\begin{youdobox}[Reading measuring cylinders]
\task Read each measuring cylinder. Give the value of one small division and the volume, with units.\skr{PC4}

\begin{center}
\begin{tabular}{ccc}
  \readcyl{4}{9}{1}{5}{6.4} & \hspace{1.2cm}\readcyl{50}{100}{10}{5}{76} & \hspace{1.2cm}\readcyl{100}{250}{50}{5}{180} \\[4pt]
  \textbf{(a)} & \hspace{1.2cm}\textbf{(b)} & \hspace{1.2cm}\textbf{(c)} \\
\end{tabular}
\end{center}
\begin{ctab}
\begin{tabular}{|SZ{1.6cm}|Z{6.2cm}|Z{6.2cm}|}
  \hline\hd & \thead{Cylinder} & \thead{Value of one small division (mL)} & \thead{Volume (mL)} \\ \hline
  & \scaf{e.g.} & \scaf{1} & \scaf{34} \\ \hline
  & (a) & & \\ \hline
  & (b) & & \\ \hline
  & (c) & & \\ \hline
\end{tabular}
\end{ctab}
\end{youdobox}

% ─── PLAN AND SAFETY / METHOD ────────────────────────────────────────────────
\exppart{Plan, safety and method}{Teal}

\textbf{Aim:} To measure mass and volume precisely, and to find out which equipment measures volume most precisely.\skr{PC1}

\subsubsection*{Equipment}
\begin{multicols}{2}
\begin{itemize}[leftmargin=1.2em, itemsep=0pt, topsep=0pt]
  \item electronic balance (reads to 0.01\,g)
  \item marble, rubber stopper, 10 paper clips
  \item 100\,mL beaker
  \item four measuring cylinders, filled by your teacher (B1--B4)
  \item 250\,mL beaker with markings
  \item 50\,mL and 100\,mL measuring cylinders
  \item wash bottle of water, paper towel
\end{itemize}
\end{multicols}

\begin{infobox}{Safety}
Keep water well away from the balance and its power cord, and wipe up any spill straight away. Carry measuring cylinders by the body, not the rim; they tip over easily, so stand them away from the bench edge. Report broken glass to your teacher.
\end{infobox}

\subsubsection*{Station A --- Mass \hfill \sk{PC6}}
\begin{enumerate}[itemsep=2pt]
  \item Tare the balance so it reads 0.00\,g.
  \item Place the first object in the middle of the pan. Record its mass in Table A.
  \item Take the object off, check the balance reads 0.00\,g again, and repeat until a different person in your group has measured it each time (three trials).
  \item Repeat for each object, then calculate each average.
\end{enumerate}

\subsubsection*{Station B --- Volume}
\begin{enumerate}[itemsep=2pt]
  \item For each cylinder B1--B4, work out the value of one small division.
  \item Read the volume at eye level, from the bottom of the meniscus. Record it in Table B. Do not pick the cylinder up.
\end{enumerate}

\subsubsection*{Station C --- The precision challenge}
\begin{enumerate}[itemsep=2pt]
  \item Put the empty 100\,mL beaker on the balance and press \textbf{TARE}.
  \item Measure 35\,mL of water using the markings on the 250\,mL beaker. Pour it into the beaker on the balance and record its mass in Table C.
  \item Empty and dry the beaker, tare again, and repeat with the 100\,mL measuring cylinder, then the 50\,mL measuring cylinder.
  \item One millilitre of water has a mass of 1\,g, so a perfect 35\,mL would read 35.00\,g. Calculate how far each result is from 35.00\,g.
\end{enumerate}

% ─── RESULTS ─────────────────────────────────────────────────────────────────
\exppart{Results}{Amber}

\subsubsection*{Table A --- The mass of four objects \hfill \sk{PA1}}
\begin{ctab}
\begin{tabular}{|SW{3.6cm}|Z{2.6cm}|Z{2.6cm}|Z{2.6cm}|Z{3.2cm}|}
  \hline\hd & \thead{Object} & \thead{Trial 1 (g)} & \thead{Trial 2 (g)} & \thead{Trial 3 (g)} & \thead{Average (g)} \\ \hline
  & \scaf{e.g. eraser} & \scaf{21.36} & \scaf{21.35} & \scaf{21.36} & \scaf{21.36} \\ \hline
  & Marble & & & & \\ \hline
  & Rubber stopper & & & & \\ \hline
  & 10 paper clips & & & & \\ \hline
  & 100\,mL beaker & & & & \\ \hline
\end{tabular}
\end{ctab}

\begin{scaffold}{Calculating an average \textbar\ PA3 \skname{PA3}}
Add the three trials and divide by 3. Round to the same number of decimal places as the readings.\quad
e.g.\ $(21.36 + 21.35 + 21.36) \div 3 = 64.07 \div 3 = 21.356\ldots \approx 21.36$\,g
\end{scaffold}

\subsubsection*{Table B --- Reading measuring cylinders}
\begin{ctab}
\begin{tabular}{|SZ{2.4cm}|Z{3.6cm}|Z{4.6cm}|Z{4.4cm}|}
  \hline\hd & \thead{Cylinder} & \thead{Size of cylinder (mL)} & \thead{Value of one small division (mL)} & \thead{Volume (mL)} \\ \hline
  & \scaf{e.g.} & \scaf{25} & \scaf{0.5} & \scaf{18.5} \\ \hline
  & B1 & & & \\ \hline
  & B2 & & & \\ \hline
  & B3 & & & \\ \hline
  & B4 & & & \\ \hline
\end{tabular}
\end{ctab}

\subsubsection*{Table C --- Measuring 35\,mL three ways}
\begin{ctab}
\begin{tabular}{|SW{5.6cm}|Z{4.4cm}|Z{5.4cm}|}
  \hline\hd & \thead{Equipment used to measure 35\,mL} & \thead{Mass of water (g)} & \thead{Difference from 35.00\,g (g)} \\ \hline
  & \scaf{e.g. markings on a conical flask} & \scaf{37.80} & \scaf{$37.80 - 35.00 = 2.80$} \\ \hline
  & Markings on the 250\,mL beaker & & \\ \hline
  & 100\,mL measuring cylinder & & \\ \hline
  & 50\,mL measuring cylinder & & \\ \hline
\end{tabular}
\end{ctab}

% ─── ANALYSIS ────────────────────────────────────────────────────────────────
\exppart{Analysis and evaluation}{Coral}

\task Explain the difference between mass and volume. Give the unit you used for each.
\answerbox{5cm}

\needspace{10\baselineskip}
\task Look at your three trials for each object in Table A. How close together were they, and what does this tell you about the precision of the balance?\skr{PC4}
\answerbox{5cm}

\needspace{14\baselineskip}
\task A group recorded these masses for a marble: 5.62\,g, 5.61\,g and 5.16\,g.\skr{PA7}
\pt{a}Which reading is the anomaly?
\answerbox{1cm}
\pt{b}Suggest what might have caused it.
\answerbox{5cm}
\pt{c}What should the group do about it before they calculate an average?
\answerbox{5cm}

\needspace{10\baselineskip}
\task Which equipment measured 35\,mL of water most precisely in Table C? Quote your results as evidence.\skr{PC4}
\answerbox{5cm}

\needspace{10\baselineskip}
\task Why is the unit written once in each column heading, instead of after every number in the table?\skr{PA1}
\answerbox{5cm}

\needspace{14\baselineskip}
\task Readings from a measuring cylinder can be wrong.\skr{EV3}
\pt{a}Describe one mistake that gives a wrong reading, and how to avoid it.
\answerbox{5cm}
\pt{b}Describe a different mistake, and how to avoid it.
\answerbox{5cm}

\needspace{8\baselineskip}
\task A beaker has a mass of 64.15\,g. With water in it, it has a mass of 112.40\,g. Calculate the mass of the water.\skr{PA3}
\workbox{3.5cm}

\begin{yr8box}
\begin{keepext}
\exttask A kitchen scale reads to the nearest 1\,g. The lab balance reads to 0.01\,g. A 10-cent coin has a mass of about 5.65\,g. Which would you use to measure the coin's mass, and why?
\answerbox{5cm}
\end{keepext}
\begin{keepext}
\exttask A student measures 7\,mL of water in a 100\,mL measuring cylinder whose small divisions are 1\,mL. Explain why this is a poor choice, and how much better a 10\,mL cylinder with 0.2\,mL divisions would be.
\answerbox{5cm}
\end{keepext}
\end{yr8box}

% ─── SKILL CHECK ─────────────────────────────────────────────────────────────
\exppart{Skill check}{NavyBlue}

\begin{skillcheck}
\scrow{PC7}{I can tare an electronic balance and record a mass with every digit and its unit.}
\scrow{PC4}{I can read a measuring cylinder at eye level from the bottom of the meniscus, and choose the most precise cylinder.}
\scrow{PC6}{I measured each object three times and recorded every reading straight away.}
\scrow{PA1}{I can record data in a table with the units in the column headings.}
\scrow{PA3}{I can calculate an average and a difference.}
\scrow{PA7}{I can spot a reading that does not fit with the others.}
\scrow{EV3}{I can name a mistake that makes a reading wrong, and how to avoid it.}
\end{skillcheck}

\task Which skill on this page will you focus on in the next experiment? Explain what you will do differently.
\answerbox{5cm}
